\section{An explicit rank cutoff for arbitrary unitary characters}
\label{sec:unitary-nonweil-cutoff}

\paragraph{Status and scope.}
This is a new author derivation from two published character estimates.
The complete proof, invoked primary sources and exact constants have
passed author, root and independent peer reviews. The peer certificate
independently reconstructs every quantitative reduction.
The new estimate concerns all ordinary non-Weil irreducible characters
of special unitary groups of sufficiently large natural dimension.
It applies to their full compatible linear normalizers, without a
self-duality hypothesis or an equality of inducing and original
character fields.  The numerical cutoff is deliberately conservative.
The theorem does not classify the characters below the cutoff.

Put
\[
 \sigma_0=10^{-29},\qquad J_0=10^{30},\qquad N_0=10^{61}.
                                                               \tag{UNW.1}
\]
Here $n$ always denotes the natural dimension of
$\operatorname{SU}_n(q)$; the degree of a complex character is denoted
by $a$.  For a finite-order complex operator $T$ on a space of
dimension $a$, write
\[
 \nu(T)=a-\max_{\lambda}\dim\ker(T-\lambda).
                                                               \tag{UNW.2}
\]
Thus $\nu$ is unchanged by multiplying $T$ by a scalar.  Ordinary
moving rank means $\operatorname{rank}(T-1)$ and is a different
quantity.

\begin{lemma}[D-UNW-01: the two quantitative published inputs]
\label{D-UNW-01}
Let $S=\operatorname{SU}_n(q)$, $n\ge3$, and let
$\chi\in\operatorname{Irr}(S)$ have degree $a$ and character level
$j$.  Then
\[
 \frac{|\chi(g)|}{a}\le
 a^{-\sigma_0\operatorname{supp}(g)/n}\qquad(g\in S),
                                                               \tag{UNW.3}
\]
where $\operatorname{supp}(g)$ is $n$ minus the largest eigenspace
dimension of $g$ on the natural module over an algebraic closure.
The level satisfies
\[
 a\ge \frac{q^{j(n-j)}}{2(q+1)},\qquad
 j\ge n/2\ \Longrightarrow\
 a>\frac{q^{n^2/4-2}}{q+1}.                                  \tag{UNW.4}
\]

If $2\le j<n/2$ and $n\ge j^2+j+1$, set
$Q=\operatorname{GU}_j(q)$ and
$R_j=\sqrt{16.52q^{j^2+j-1}}$.
There is an $\alpha\in\operatorname{Irr}(Q)$ and an extension of
$\chi$ to $\operatorname{GU}_n(q)$, of the same degree, such that
for every noncentral $g\in S$, with
$k=n-\operatorname{supp}(g)\le n-1$,
\[
\begin{split}
 a&\ge\frac{\alpha(1)}{|Q|}
 \left(q^{nj}-|Q|q^{n(j-1)}-|Q|R_jq^{n(j-2)}\right),\\
 |\chi(g)|&\le\frac{\alpha(1)}{|Q|}
 \left(q^{kj}+|Q|q^{(n-1)(j-1)+1}
                  +|Q|R_jq^{n(j-2)}\right),\\
 |Q|&\le1.5q^{j^2}.
\end{split}                                                    \tag{UNW.5}
\]
These estimates include $q=2$.
\end{lemma}

\begin{proof}
We specify the known inputs and their constants rather than reprove
the published character theorems.
Larsen--Tiep prove (UNW.3) with an explicit larger constant
\[
 \sigma=\frac1{6507\cdot2^{21}\cdot10^{18}}>7\cdot10^{-29}.
                                                               \tag{UNW.6}
\]
This is Theorem~5.5 and the last displayed computation of its
complete proof in \cite[accepted-author version, pp.~25--26]{UNW:LT2024}.
The constants supplied in their Theorem~4.4 are
$C=2^{14}10^{12}$ and $\gamma=(2^{13}10^{12})^{-1}$, and the proof
takes
$\sigma=\min\{[241(9C)^{3/2}]^{-1},\gamma/82\}$.
Consequently the smaller rational constant $\sigma_0$ in (UNW.1)
is valid.  The proof read here is the accepted author version;
the journal publication and DOI are verified separately.  No claim
is made to have retrieved its typeset journal PDF.

We use the level defined for $S$ in
\cite[Definition~2 and Lemma~7.1, pp.~51--52]{UNW:GLT2020}:
it is the level of any irreducible character of
$\operatorname{GU}_n(q)$ lying above $\chi$.
Lemma~7.1 proves independence of the choice.  Equations (UNW.4)
are precisely Theorem~7.3(i)--(ii) there, with
$\sigma_-=1/[2(q+1)]$.

For $j<n/2$, Corollary~8.6(i), including its proof on pp.~67--68,
shows that a character above $\chi$ restricts irreducibly to $S$.
Its determinant twists contain a unique character of true level
$j$.  The determinant twists are trivial on $S$.
Theorem~1.1(iii), proved in full in \S8.3 on pp.~65--67, identifies
this true-level character with $D_\alpha^\circ$ for some
$\alpha\in\operatorname{Irr}(\operatorname{GU}_j(q))$ whenever
$j^2+j\le n-1$.  This condition is sufficient in even as well as
odd characteristic.  The exact estimate (8.18), p.~66, gives
(UNW.5); the bound on $|Q|$ is recorded immediately below it.
The tail estimate used in (8.18) is (8.17), and the complete
subsequent argument proves that these characters are irreducible,
pairwise distinct and exhaust true level $j$.
We use this exact estimate, rather than its weaker numerical
corollary with coefficient $2.43$.
\end{proof}

\begin{theorem}[D-UNW-02: a uniform non-Weil character ratio]
\label{D-UNW-02}
Let $n\ge N_0$ and $q\ge2$, and let
$\chi\in\operatorname{Irr}(\operatorname{SU}_n(q))$ have level
$j\ge2$ and degree $a$.  Then every noncentral $g$ satisfies
\[
                         |\chi(g)|\le\frac4{15}a.            \tag{UNW.7}
\]
Moreover $a>q^n>30$.  Level zero is the trivial representation;
level-one characters are the ordinary unitary Weil characters
and have degree at most $q^n$.
\end{theorem}

\begin{proof}
First record a convenient degree consequence of (UNW.4).  Since
$2(q+1)\le q^3$ and $q+1\le q^2$, for $n\ge12$ we have
\[
 a\ge q^{nj/4}\quad(1\le j\le n/2),\qquad
 a\ge q^{n^2/8}\quad(j\ge n/2).                             \tag{UNW.8}
\]
For the first assertion, $j(n-j)-3\ge nj/2-3\ge nj/4$
once $nj\ge12$, as is automatic here.  For the second,
$n^2/4-4\ge n^2/8$ for $n\ge8$.

If $j>J_0$, put $j'=\min\{j,n/2\}$.  Because $n\ge N_0>2J_0$,
we have $j'\ge J_0$.  Every noncentral natural operator has support
at least one.  Equations (UNW.3) and (UNW.8) therefore give
\[
 \frac{|\chi(g)|}{a}
 \le q^{-\sigma_0 j'/4}
 \le 2^{-\sigma_0J_0/4}
 =2^{-5/2}<\frac14<\frac4{15}.                             \tag{UNW.9}
\]
The high-level argument does not assume that $\chi$ extends to
the general unitary group.

It remains to consider $2\le j\le J_0$.  Our constants satisfy
$N_0\ge J_0^2+J_0+8$, so (UNW.5) applies and $n-j^2\ge8$.
Use the deliberately weaker estimates
\[
                    |Q|\le2q^{j^2},\qquad R_j\le5q^{j^2}.
                                                               \tag{UNW.10}
\]
The second follows from $16.52<25$ and
$j^2+j-1\le2j^2$.  Define
\[
 u=q^{j^2-n}\le2^{-8},\qquad d=2u+10u^2.
\]
After dividing the lower bound in (UNW.5) by
$\alpha(1)q^{nj}/|Q|$, its right side is at least $1-d$.
The normalized numerator in the upper bound is at most
\[
 q^{-j}+2q^{j^2-n-j+2}+10q^{2j^2-2n}
 \le q^{-j}+d\le\frac14+d.
\]
Here $j\ge2$ is essential.  Exact rational arithmetic gives
\[
 d\le\frac2{256}+\frac{10}{256^2}
      =\frac{261}{32768}<\frac1{100},\qquad
 \frac{1/4+1/100}{1-1/100}=\frac{26}{99}<\frac4{15}.
                                                               \tag{UNW.11}
\]
The denominator is positive, so (UNW.7) follows for every low
level as well, including $q=2,j=2$.

For completeness the degree separation is uniform.  If
$2\le j\le n/2$, monotonicity of $j(n-j)$ on this interval and
(UNW.4) give $a\ge q^{2n-7}>q^n$ for $n\ge8$.
If $j\ge n/2$, they give $a>q^{n^2/4-4}>q^n$ for $n\ge8$.
Theorem~7.3(i) gives the upper bound $q^n$ at level one.
Definition~2 and Lemma~7.1 identify level zero with the trivial
character.  The identification at level one is also explicitly
recorded at the end of \S8.3, p.~67, of \cite{UNW:GLT2020}.
\end{proof}

\begin{theorem}[D-UNW-03: full normalizers and the finite-cover cutoff]
\label{D-UNW-03}
Let a finite group $J\subset\operatorname{GL}(V)$ normalize an
irreducibly acting central product of quasisimple groups
$E_1\cdots E_t$.  Suppose that, on the associated tensor factors,
each $E_i$ acts through a nontrivial ordinary level-at-least-two
representation of $\operatorname{SU}_{n_i}(q_i)$ with
$n_i\ge N_0$.  Factor permutations are allowed, including all
permutations compatible with the represented central quotients.
Then every nonscalar $T\in J$ satisfies
\[
                         \nu(T)\ge\frac{11}{60}\dim V.
                                                               \tag{UNW.12}
\]
Such a representation cannot be an inducing block of an actual
finite-cover Weil sixfold source with period rank nine.

In particular, in the primitive factorwise tensor-indecomposable
normal-core situation of D-BI-19, if the simple factors are
$\operatorname{PSU}_n(q)$, a dominant source must satisfy
\[
                              n<10^{61}.                    \tag{UNW.13}
\]
This conclusion combines the new non-Weil estimate with the
already proved unitary Weil theorem D-BI-24.  In the nonsplit
main-source scope of D-BI-20 the remaining general unitary gate
therefore has both $n<10^{61}$ and $q\le27$.
\end{theorem}

\begin{proof}
We first prove a full-normalizer estimate for one factor of degree
$a$.  Average a positive Hermitian form over its finite image.
If $T$ is a nonscalar normalizer, Schur's lemma gives an element
$x$ of the image such that $b=TxT^{-1}x^{-1}\ne1$.
Choose an eigenvalue $\lambda$ of $T$ with largest eigenspace.
The rank identity for a commutator gives
\[
 \operatorname{rank}(b-1)
 =\operatorname{rank}(Tx-xT)
 \le2\operatorname{rank}(T-\lambda)=2\nu(T).                \tag{UNW.14}
\]
The operator $b$ lies in the finite core image.  If it is a
nonidentity scalar, its ordinary moving rank is $a$.
Otherwise it is represented by a noncentral element of the
special unitary group and (UNW.7) applies.  Its eigenvalues have
absolute value one, so
\[
 2\operatorname{rank}(b-1)
 \ge a-\operatorname{Re}\operatorname{Tr}(b)
 \ge(1-4/15)a=11a/15.
\]
Together with (UNW.14) this proves $\nu(T)\ge11a/60$.
No assumption on which outer automorphism $T$ implements has
entered the proof.  Represented central quotients also cause no
problem: lift the character to $\operatorname{SU}_n(q)$, and
apply the estimate to any preimage of the commutator.

If an element fixes all tensor factors, it is a tensor product
of the corresponding intertwiners.  An eigenspace proportion of
a tensor product is at most the largest eigenspace proportion
on any chosen factor: this follows by summing the products of
the eigenvalue multiplicities.  If the whole operator is
nonscalar, some factor is nonscalar, giving (UNW.12).

For clarity the factor-permutation estimate does not require
any character-field assumption.  On a cycle of length
$\ell\ge2$ with factor degree $a\ge2$, changes of tensor bases
reduce the weighted cyclic permutation to a cyclic shift with
one return operator.  This return operator is diagonalizable:
a power is scalar because the original operator has finite
projective order.  In its eigenbasis, tensor words are permuted
in cyclic orbits, each of which contributes at most one copy
of a prescribed eigenvalue.  There are $a$ constant words and
every other orbit has length at least two.  The number of orbits
is at most
\[
 a+(a^\ell-a)/2\le3a^\ell/4.
\]
Hence such a cycle has $\nu\ge a^\ell/4$; convolution with the
other factors preserves the proportion.  Since $1/4>11/60$,
this proves (UNW.12) in every case.

We spell out the geometric constraint used to convert this
estimate into an obstruction.  Let $e=\dim V$ be an inducing
degree, let $X$ be the intermediate base curve of genus $h$,
and write $r$ for its nontrivial inertia points.  Polarized
Shapiro and the finite-cover Hurwitz factorization, already
proved for this source problem, identify the two chosen Hodge
multiplicities with $(3,3)$.  For the irreducible nontrivial
coefficient representation the Euler identity is
\[
 6=2e(h-1)+\sum_{i=1}^r\operatorname{rank}(g_i-1).
                                                               \tag{UNW.15}
\]
The additional finite-cover data over a pointed $X$ are discrete,
so the period map factors through a finite cover of its moduli.
For $h\ge2$ the right side is already greater than six, since
$e>30$.  For $h=1$, the sum of moving ranks is six, hence
$r\le6$ and the pointed-curve moduli dimension is at most six
(or one if $r=0$).  Such a source cannot have period rank nine.
Therefore $h=0$, $r\ge12$, and
\[
                  \sum_i\operatorname{rank}(g_i-1)=2e+6.
                                                               \tag{UNW.16}
\]
Every nonidentity scalar inertia has moving rank $e$, and every
nonscalar inertia has moving rank at least $\nu(g_i)\ge11e/60$.
Thus
\[
               2e+6\ge12\cdot\frac{11e}{60}
                  =\frac{11e}{5},\qquad e\le30.
                                                               \tag{UNW.17}
\]
Each factor already has degree greater than thirty by
D-UNW-02, a contradiction.

Finally consider the transitive core of D-BI-19.  At these large
natural dimensions there is no exceptional multiplier; the
universal perfect cover of $\operatorname{PSU}_n(q)$ is
$\operatorname{SU}_n(q)$, and all its quasisimple covers are
central quotients.  The precise finite-exception boundary is
the published list in \cite[Chapter~5, Table~1, p.~35]{BI:HHM},
with the ordinary-cover distinction on p.~4.  Thus every factor
character lifts to the special unitary group.  Transitive
conjugation makes all factor degrees equal.  By the degree
separation at the end of D-UNW-02, either they all have level
one or they all have level at least two.  The latter case was
just excluded.  The former is excluded by D-BI-24; none of
its small unitary exceptions has natural dimension at least
$N_0$.  This proves (UNW.13).  The additional field cutoff is
exactly D-BI-20, with its nonsplit and primitive-core hypotheses
retained.
\end{proof}

\paragraph{What the finite reduction does and does not decide.}
The exact integer $10^{61}$ is an explicit upper cutoff for the
remaining natural unitary rank in the stated main-source problem.
Together with the prior field bound it leaves a finite collection
of finite groups, each with finitely many ordinary characters.
It does not supply their tables, classify all lower-rank tuples,
or prove existence of any remaining source.  Nor does (UNW.12)
exclude an arbitrary mixed core containing additional uncovered
small-rank factors: its mixed statement requires every factor
to satisfy the displayed large-rank non-Weil hypotheses.

\paragraph{Certificate.}
The standalone standard-library program
\texttt{unitary\_nonweil/check\_constants.py} checks all numerical
constants using exact integers and rational numbers, together
with the polynomial identities underlying the two degree regimes.
It does not enumerate ranks up to $10^{61}$; universal monotonicity
and the two-regime proof are given above.  The retained source
retrieval manifest records the exact full-text hashes and the
unsuccessful attempts to obtain a typeset LT journal PDF.
